\documentclass[12pt]{article}
\usepackage[pdftex]{graphicx}
\usepackage[T1]{fontenc}
\usepackage{textcomp}
\usepackage{hyperref}
\hypersetup{
    colorlinks,
    citecolor=black,
    filecolor=black,
    linkcolor=black,
    urlcolor=cyan
}
\usepackage[text={16cm,24cm}]{geometry}
\usepackage{pgffor}
\usepackage{filecontents}% Used so that the external files can be placed in this file
\setlength{\parskip}{\baselineskip}
\usepackage{color}
\definecolor{shadowgrey}{rgb}{0.18,0.2,0.21}
\usepackage{listings}
\lstset{ %
  frame=shadowbox,
  rulesepcolor=\color{shadowgrey},
  basicstyle=\footnotesize,
  commentstyle=\footnotesize
}

\newcommand{\todo}[1]{\colorbox{yellow}{\textbf{TODO: #1} fill this in!}}
\newcommand{\linuxtilde}{\raise.20ex\hbox{$\scriptstyle\mathtt{\sim}$}}
\newcommand{\link}[1]{\hyperref[sec:#1]{\textbf{section~\ref*{sec:#1}}}}

\title{\textbf{EtherCAT Dextrous Hand} \\
- User Manual -}
\author{Shadow Software team - software@shadowrobot.com}

\makeindex

\begin{document}
\begin{titlepage}

\maketitle
\vspace{5cm}
\begin{center}
\includegraphics{images/logo-shadowDB.png}
\end{center}
\end{titlepage}

\tableofcontents
\newpage

\section{Overview}
\label{sec:overview}

\par This user manual will get you started with using our EtherCAT hand or simulated robots. We'll guide you through installing the code in the \link{install}, show you some basic commands and code samples in \link{running-our-code}, and explain where to find things in \link{navigate}. This guide ends with some pointers to essential documentation you should read (\link{where-to-go}), and some troubleshooting tips (\link{what-do-if}).

\par If you have any questions, don't hesitate to contact us at: \texttt{software@shadowrobot.com}.

\newpage

\section{Installing our Software}
\label{sec:install}
\par For a typical user, we really recommend installing our software directly from the apt repository (see section \ref{sec:install-apt}). This will ensure that you keep the latest released software on your computer.

\par We assume you've already installed ROS electric, as detailed on \href{http://www.ros.org/wiki/electric/Installation/Ubuntu}{the ROS Wiki}.

\par However if you plan to modify some parts of our code, you can use an overlay to install some of the sources from our publicly available Launchpad repository on top of your existing installation.

\subsection{From our apt repository (recommended)}
\label{sec:install-apt}
\todo{Mark}

\subsubsection{Install a sr\_config overlay for the EtherCAT Hand}
\label{sec:install-cfg}
\par If you have an EtherCAT Hand, you should create an overlay for the \texttt{sr\_config} stack: this stack contains all the different configuration files which can be modified to reflect your specific system. You shouldn't use the one installed by default as they are defined for the simulated hand.

\par To create this overlay, just follow the \link{install-src}.

\subsection{From sources}
\label{sec:install-src}
\par If you need to modify one of the packages for your own use, you'll want to first install the released version from the apt repository as explained in the \link{install-apt}. This will ensure that you have all the code you need. Once you've identified in which stack the package you want to modify is, you can create an overlay\footnote{\textbf{Overlay:} use another version of a stack than the default installed one.} to use the source version.

\par To create an overlay, we'll use the \href{http://ros.org/wiki/rosinstall}{rosinstall} utility. In this example, we'll overlay the \href{http://launchpad.net/sr-config}{sr\_config} stack, using trunk version of the code hosted on launchpad using \href{http://bazaar.canonical.com}{bzr}: \texttt{lp:sr-config}.

\begin{itemize}
\item First we need to install the rosinstall utility if you haven't done so already:
  \begin{lstlisting}[language=Bash]
> sudo apt-get install python-setuptools bzr
> sudo easy_install -U rosinstall vcstools
  \end{lstlisting}

\item We can now create a rosinstall file that will contain the information about the stacks we want to overlay (in our case we're overlaying the sr\_config stack only):
  \begin{lstlisting}[title={\textbf{shadow\_robot.rosinstall}}, language=Python]
- bzr:                       # VCS used (bzr in our case)
  uri: 'lp:sr-config'        # public uri for the code
  local-name: sr_config      # name under which we install the stack
  \end{lstlisting}

\item We're now going to create a workspace (in the \texttt{\linuxtilde/workspace} directory), based on your currently installed ROS system (we're assuming you're using ROS electric), and use the rosinstall file we created to download the code into it.
  \begin{lstlisting}[escapeinside='', language=Bash]
> rosinstall '\linuxtilde'/workspace /opt/ros/electric/setup.bash
> rosinstall '\linuxtilde'/workspace shadow_robot.rosinstall
  \end{lstlisting}

\item To use the newly created environment, simply source the setup.bash created in the previous step (to have it sourced automatically when you open a new terminal, simply add it to your \texttt{\linuxtilde/.bashrc} as show on the second line).
  \begin{lstlisting}[escapeinside='', language=Bash]
> source '\linuxtilde'/workspace/setup.bash
> echo source '\linuxtilde'/workspace/setup.bash >> '\linuxtilde'/.bashrc
  \end{lstlisting}

\item Finally, you can build the whole stack in one go using the \texttt{rosmake} command:
  \begin{lstlisting}[escapeinside='', language=Bash]
> rosmake --rosdep-install --rosdep-yes sr_config
  \end{lstlisting}

\end{itemize}

\subsubsection{Available rosinstall files}
\par If you want to install all our stacks in one go, you can use the \texttt{rosinstall} files available \href{http://bazaar.launchpad.net/~shadowrobot/sr-build-tools/trunk/view/head:/data/shadow_robot.rosinstall}{here}.

\todo{Mark} May be we should have a simpler shadow\_robot.rosinstall which doesn't contain the unstable stacks (sr-demo / denso / object-manip / sr-manip)?

\par Once you've downloaded the code using the \texttt{rosinstall} commands described previously, you can build all our stacks:
  \begin{lstlisting}[escapeinside='', language=Bash]
> rosmake --rosdep-install --rosdep-yes shadow_robot \
  shadow_robot_ethercat sr_contrib sr_visualization \
  sr_teleop
  \end{lstlisting}

\subsection{Import our custom Environment Variables}
\label{sec:import-vars}
\par As described in \link{spec-system-param}, we're using a file to define some useful environment variables. To make sure the variables are imported in your setup, you need to source it in your \texttt{\linuxtilde/.bashrc}, after sourcing the ros setup.bash. It should look like this:
  \begin{lstlisting}[title={\textbf{\linuxtilde/.bashrc}}, upquote=true]
#source the ros setup.bash
/opt/ros/electric/setup.bash
#or the one created in your workspace
# if you used an overlay
#source ~/workspace/setup.bash

#now source our additional environment variables
source `rosstack find sr_config`/bashrc/env_variables.bashrc
  \end{lstlisting}

\newpage

\section{Navigating our code}
\label{sec:navigate}

\subsection{Important Packages}
\label{sec:important-packages}



\subsection{Where to find the settings}
\label{sec:where-find-settings}

\subsubsection{Specifying system parameters}
\label{sec:spec-system-param}

\subsubsection{Specifying default controllers settings}
\label{sec:spec-defa-contr}

\newpage

\section{Running our code}
\label{sec:running-our-code}

\subsection{Starting the real hand}
\label{sec:starting-real-hand}

\subsubsection{Setting up the Ethernet port}
\par The first time you plug an EtherCAT Hand, you probably need to check on which Ethernet port the hand is connected (you probably have 2 network ports in your computer, one of them would typically be connected to the network, the other one to the hand).

\par An easy way to see which port the hand is using is to unplug the hand network cable from your computer and then plug it back in. You can now type \texttt{dmesg} and you should see a line saying which port has just been plugged in, for example, in the lines below, you see that \texttt{eth1} is used by the hand:
  \begin{lstlisting}[escapeinside='', language=Bash]
[16661.901920] e1000e: eth1 NIC Link is Down
[16666.460518] e1000e: eth1 NIC Link is Up 1000 Mbps Full Duplex
  \end{lstlisting}

\par Now that we know which Ethernet port is used by the hand (\texttt{eth1} in our example), we can edit the interface in the system parameters, as shown in \link{spec-system-param}:
\begin{itemize}
\item Make sure you've created your overlay for the \texttt{sr\_config} stack as described in the \link{install-cfg}.
\item Let's open the configuration file:
  \begin{lstlisting}[escapeinside='', language=Bash]
> roscd sr_config
> cd bashrc
> gedit env_variables.bashrc
  \end{lstlisting}
\item Set the \texttt{ETHERCAT\_PORT} to the correct value (it is \texttt{eth1} by default)
\item Save and restart your terminal to make sure the new options have been taken into account.
\end{itemize}

\subsubsection{Starting the driver}
\par You can now start the driver. To be able to run the driver, you need to elevate the permissions to root access:
  \begin{lstlisting}[escapeinside='', language=Bash]
#elevate permissions
> sudo -s
> roslaunch sr_edc_launch sr_edc.launch
  \end{lstlisting}

\subsection{Starting the simulated robot}
\label{sec:start-simul-hand}
\par To start the simulated robot you can simply run:
  \begin{lstlisting}[escapeinside='', language=Bash]
> roslaunch sr_hand gazebo_hand_motor.launch
  \end{lstlisting}

\par Or if you want a simulated Shadow Arm as well as our hand:
  \begin{lstlisting}[escapeinside='', language=Bash]
> roslaunch sr_hand gazebo_arm_and_hand_motor.launch
  \end{lstlisting}

\subsection{Checking it works}
\label{sec:checking-it-works}

\newpage

\section{Where to go next}
\label{sec:where-to-go}

\newpage

\section{What to do if it doesn't work}
\label{sec:what-do-if}

\newpage

\section{Changelog}
\label{sec:changelog}

\newcommand{\filename}[1]{changelogs/file#1}
\newcommand{\buildchangelog}
{%
  \foreach \i in {0, ...,99}%
  {%
    \expandafter\IfFileExists\expandafter{\expandafter\filename\expandafter\i}%
    {%
      \expandafter\input\expandafter{\filename\i}%
    }%
    {}%
  }%
}
\buildchangelog

\end{document}

